\uuid{7m8Y}
\exo7id{3477}
\titre{exo7 3477}
\auteur{quercia}
\organisation{exo7}
\datecreate{2010-03-10}
\isIndication{false}
\isCorrection{false}
\chapitre{Déterminant, système linéaire}
\sousChapitre{Système linéaire, rang}
\cpge{mpsi-mp2i.changements-bases.b, mpsi-mp2i.matrices-applications-lineaires.b, pcsi.matrices-applications-lineaires.b, ptsi.matrices-applications-lineaires.b, mpsi-mp2i.matrices-applications-lineaires.a, pcsi.matrices-applications-lineaires.a, ptsi.matrices-applications-lineaires.a}
\module{Algèbre}
\niveau{L2}
\difficulte{}

\contenu{
\question{
Soient $A \in \mathcal{M}_{n,p}(K)$ et $B \in \mathcal{M}_{q,r}(K)$. Chercher le rang des
applications :
$$
   f : {\mathcal{M}_{p,q}(K)} \to {\mathcal{M}_{n,q}(K)}, X \mapsto {AX}\qquad
   g : {\mathcal{M}_{p,q}(K)} \to {\mathcal{M}_{p,r}(K)} X \mapsto{XB}  $$
$$  h  : {\mathcal{M}_{p,q}(K)} \to {\mathcal{M}_{n,r}(K)} X \mapsto{AXB}
$$
(on transformera $A$ et $B$ en matrices canoniques équivalentes)
}
}
